\documentclass[11pt]{article}
\usepackage{amsmath,amssymb,amsthm,booktabs,mathtools}
\usepackage[margin=1in]{geometry}
\usepackage{microtype}
\usepackage{url}
\newtheorem{theorem}{Theorem}[section]
\newtheorem{lemma}[theorem]{Lemma}
\newtheorem{remark}[theorem]{Remark}
\newcommand{\dd}{\,\mathrm d}
\newcommand{\PhiH}{\Phi_H}
\title{Atlas 126: A Triangle--$C_4$ Vertex-Amalgam Supporting Plane}
\author{}
\date{}
\begin{document}
\maketitle

\begin{abstract}
Let $H$ be a triangle and a $4$-cycle sharing exactly one vertex.  We prove
that every graphon of edge density $p\in[1/2,1]$ satisfies
\[
 t(H,W)\ge p^2(2p-1)\bigl(p^3+(1-p)^3\bigr).
\]
This is the full chromatic interval and the exact target for Graph Atlas 126.
Rooting at the common vertex factors the density into a rooted triangle and a
rooted $C_4$.  A two-piece weighted Cauchy--Schwarz projection then reduces
the graphon problem to a scalar function of the degree, the one-step degree
average, and the rooted triangle density.  Its sharp supporting plane uses
Reiher's triangle-density theorem on $1/2\le p\le2/3$ and is elementary above
$2/3$.  The remaining rational polynomial inequalities are proved by an
exact Bernstein subdivision certificate, including exact local analyses of
all zero-coefficient boxes.
\end{abstract}

\section{The graph and its chromatic target}

Let $(\Omega,\mu)$ be an arbitrary probability space and let
$W\colon\Omega^2\to[0,1]$ be symmetric and measurable.  All equalities and
inequalities below hold after changing functions on null sets if necessary.
Every integrand is nonnegative and bounded, so Tonelli's theorem justifies
all changes in the order of integration.

The Atlas-labelled edge set of Graph Atlas 126 is
\[
 \{02,05,13,14,23,34,35\}.
\]
It has graph6 string \texttt{ERUO} and degree sequence
$(4,2,2,2,2,2)$.  Its triangle is $1$--$3$--$4$--$1$, and its $4$-cycle is
$0$--$2$--$3$--$5$--$0$; these subgraphs share only vertex $3$.
Thus it is isomorphic to a triangle and a $C_4$ with one vertex identified.

If $r$ colours are available, first colour the common vertex.  The remaining
two triangle vertices have $(r-1)(r-2)$ choices.  The three remaining cycle
vertices have $(r-1)(r^2-3r+3)$ choices.  Hence
\begin{equation}
 P_H(r)=r(r-2)(r-1)^2(r^2-3r+3),\qquad \chi(H)=3.
 \label{eq:chromatic}
\end{equation}
The chromatic target is
\begin{align}
 \PhiH(p)
 &=(1-p)^6P_H\!\left(\frac1{1-p}\right)\notag\\
 &=p^2(2p-1)\bigl(p^3+(1-p)^3\bigr).
 \label{eq:target}
\end{align}
The required interval is therefore $p\in[1/2,1]$.

\section{Rooting and the Hilbert-space projection}

Write
\[
 p=\int_{\Omega^2}W(x,y)\dd\mu(x)\dd\mu(y),\qquad
 d(x)=\int_\Omega W(x,y)\dd\mu(y),
\]
and define
\begin{align*}
 a(x)&=(T_Wd)(x)=\int_\Omega W(x,y)d(y)\dd\mu(y),\\
 t(x)&=\int_{\Omega^2}W(x,y)W(x,z)W(y,z)
       \dd\mu(y)\dd\mu(z).
\end{align*}
Thus $t(x)$ is the triangle density rooted at $x$.  For fixed $x$, put
\[
 b_x(z)=\int_\Omega W(x,y)W(y,z)\dd\mu(y).
\]
The $C_4$ density rooted at $x$ is
\[
 r_4(x)=\int_\Omega b_x(z)^2\dd\mu(z).
\]
Moreover,
\[
 \int b_x=a(x),\qquad
 \int W(x,z)b_x(z)\dd\mu(z)=t(x).
\]
Splitting $r_4(x)$ with weights $W(x,z)$ and $1-W(x,z)$ and applying
Cauchy--Schwarz on each part gives, for $0<d(x)<1$,
\begin{equation}
 r_4(x)\ge \frac{t(x)^2}{d(x)}
 +\frac{(a(x)-t(x))^2}{1-d(x)}.
 \label{eq:projection}
\end{equation}
Conditioning on the common vertex makes the triangle and $C_4$ factors
independent.  Consequently
\begin{equation}
 t(H,W)=\int_\Omega t(x)r_4(x)\dd\mu(x)
 \ge\int_\Omega F_p(d(x),a(x),t(x))\dd\mu(x),
 \label{eq:rooted-reduction}
\end{equation}
where
\begin{equation}
 F_p(d,a,t)=\frac{t^3}{d}+\frac{t(a-t)^2}{1-d}.
 \label{eq:F}
\end{equation}
At $d=0$ and $d=1$ this is interpreted by the actual rooted limits described
in Section~\ref{sec:endpoints}; no assertion about arbitrary scalar triples
at those endpoints is being made.

\section{The exact scalar feasible region}

The rooted quantities obey
\begin{equation}
 0\le d\le1,\qquad
 \max\{0,d+p-1\}\le a\le d,\qquad
 \max\{0,2a-p\}\le t\le\min\{a,d^2\}.
 \label{eq:feasible}
\end{equation}
Indeed, $a\le d$ follows from $d(y)\le1$.  Expanding
$\int(1-W(x,y))(1-d(y))\dd\mu(y)\ge0$ gives $a\ge d+p-1$.
The inequalities $t\le a$ and $t\le d^2$ follow by deleting, respectively,
one and two edge factors from the rooted triangle integral.  Finally,
\[
 W(x,y)W(x,z)\ge W(x,y)+W(x,z)-1;
\]
multiplication by $W(y,z)$ and integration give $t\ge2a-p$.

The three global moments needed below are
\begin{equation}
 \int d=p,\qquad \int a=\int d^2,\qquad
 \int t=t(K_3,W).
 \label{eq:moments}
\end{equation}
The middle identity follows directly from symmetry and Tonelli:
$\int_xa(x)=\int_yd(y)^2$.

For $p\in[1/2,2/3]$, introduce the unique $x\in[0,1]$ such that
\begin{equation}
 p=\frac{3+2x-x^2}{6},\qquad
 g(p)=\frac{x(3-x)^2}{18}.
 \label{eq:triangle-profile}
\end{equation}
The $K_3$ case of Reiher's clique-density theorem gives
\begin{equation}
 t(K_3,W)\ge g(p).
 \label{eq:reiher}
\end{equation}
This is the only prior non-elementary graph-density theorem used here.

\section{The supporting planes}

Put $q=1-p$, $c=2p-1$, and
\[
 C(p)=p^2c(p^3+q^3),\qquad D_0=p^2+q^2,\qquad h=1-\frac p2.
\]
Define the low coefficients
\begin{equation}
 \gamma_L=\frac1{16}+\frac c2+c^2,\qquad
 \rho=\frac{p^2}{D_0},\qquad
 \beta_L=2\gamma_L\rho.
 \label{eq:low-coefficients}
\end{equation}

On region $A$, $0\le x\le13/20$, set
\begin{equation}
 \gamma=\gamma_L,\qquad \beta=\beta_L,\qquad
 \lambda=
 \frac{C+\beta_L(\rho-p)+\gamma_L(p/2-\rho^2)}{g}
 +\frac{x}{1000}.
 \label{eq:region-a}
\end{equation}
The quotient has a removable singularity at $x=0$; its prescribed value is
$\lambda(0)=1/16$.  The small final term vanishes at the threshold and
strictly separates one otherwise isolated boundary contact.

On region $C$, $13/20\le x\le1$, use
\begin{equation}
 x=\frac{13+7w^2}{20},\qquad 0\le w\le1,\qquad
 \varphi=2w-w^2.
 \label{eq:transition-parameter}
\end{equation}
Set
\begin{align}
 \gamma_H&=2p(9p^2-10p+3),\notag\\
 \beta_H&=p(30p^3-29p^2+6p+1),\label{eq:high-coefficients}\\
 \gamma&=(1-\varphi)\gamma_L+\varphi\gamma_H,\notag\\
 \beta&=(1-\varphi)\beta_L+\varphi\beta_H,\notag\\
 \lambda&=\frac{C+\beta(h-p)+\gamma(p/2-h^2)}{g}.
 \label{eq:region-c}
\end{align}
The endpoint-flat interpolation is important at $p=2/3$.  At $w=1$,
\begin{equation}
 \beta=\frac23,\qquad \gamma=\frac49,\qquad \lambda=0.
 \label{eq:join-coefficients}
\end{equation}
On region $H$, $2/3\le p\le1$, use \eqref{eq:high-coefficients} directly
and set $\lambda=0$.

The exact univariate Bernstein checks in the companion verifier prove
$\gamma\ge0$ and $\lambda\ge0$.  They also prove the endpoint scalar gaps
used in Section~\ref{sec:endpoints}.  The Bernstein degrees and numbers of
zero coefficients for $(\gamma,\lambda,G_0,G_1)$ are
\begin{center}
\begin{tabular}{c@{\qquad}cccc}
\toprule
region&$\gamma$&$\lambda$&$G_0$&$G_1$\\
\midrule
$A$&$(4,0)$&$(16,0)$&$(14,0)$&$(17,0)$\\
$C$&$(14,0)$&$(30,2)$&$(30,0)$&$(34,0)$\\
$H$&$(3,0)$&identically zero&$(5,0)$&$(5,2)$\\
\bottomrule
\end{tabular}
\end{center}
where
\begin{align*}
 G_0&=\beta p+\lambda g-C,\\
 G_1&=p^3-C-\beta(1-p)-\gamma(p-1)-\lambda(p-g).
\end{align*}

\begin{lemma}[Scalar supporting plane]
\label{lem:scalar}
Let $p\in[1/2,1)$ and let $(d,a,t)$ satisfy \eqref{eq:feasible} with
$0<d<1$.  With the coefficients from the appropriate region above,
\begin{equation}
 F_p(d,a,t)\ge
 C(p)+\beta(d-p)+\gamma(a-d^2)+\lambda(t-g),
 \label{eq:supporting-plane}
\end{equation}
where $g=0$ and $\lambda=0$ in region $H$.
The same inequality holds at the actual rooted endpoint triples and at
$t=0$ under the conventions of Section~\ref{sec:endpoints}.
\end{lemma}

\section{Convex face reduction}

Fix $d$ and $t>0$.  After moving the right side of
\eqref{eq:supporting-plane} to the left, the part depending on $a$ is
\[
 \frac{t(a-t)^2}{1-d}-\gamma a.
\]
It is convex and has its unconstrained minimum at
\begin{equation}
 a_*=t+\frac{\gamma(1-d)}{2t}>t.
 \label{eq:astar}
\end{equation}
For fixed $d,t$, \eqref{eq:feasible} says
\[
 \max\{t,d-q\}\le a\le\min\left\{d,\frac{t+p}{2}\right\}.
\]
Since $a_*>t$, a minimum cannot occur only on the face $a=t$; if that face
is a singleton it is already one of the upper faces.  It is enough to check
\begin{equation}
 a=d-q,\qquad a=d,\qquad a=\frac{t+p}{2},\qquad a=a_*.
 \label{eq:four-candidates}
\end{equation}

For the interior candidate, feasibility of \eqref{eq:astar} is expressed by
\begin{align}
 K_1&=pt-t^2-\gamma(1-d)\ge0,\notag\\
 K_2&=2t(d-t)-\gamma(1-d)\ge0,\notag\\
 K_3&=2t(t-d+q)+\gamma(1-d)\ge0,\label{eq:interior-constraints}\\
 K_4&=d^2-t\ge0.
\end{align}
The redundant inequality $2t-\gamma\ge0$ follows from
\[
 2K_1+K_3=(1-d)(2t-\gamma).
\]
Substitution of \eqref{eq:astar} into the residual and multiplication by
$4dt>0$ gives the verifier target
\begin{equation}
 4dt\left(
 \frac{t^3}{d}+\gamma d^2-\beta(d-p)-C+\lambda g
 -(\gamma+\lambda)t-\frac{\gamma^2(1-d)}{4t}
 \right)\ge0.
 \label{eq:interior-polynomial}
\end{equation}
Its denominator depends only on the region parameter and has a strictly
positive Bernstein expansion.  The certificate takes $t=pv$ on the unit
cube in the region parameter, $d$, and $v$, pruning a box only when one of
\eqref{eq:interior-constraints} or $2t-\gamma$ is strictly negative there.

For completeness, the low-density boundary cubes are as follows.  Here
$u,v\in[0,1]$ and $h=1-p/2$:
\begin{center}
\begin{tabular}{ccl}
\toprule
face&range substitution&$t$\\
\midrule
$a=d-q$, I&$d=q+pu/2$&$(d-q)v$\\
$a=d-q$, II&$d=h+pu/2$&$2d+p-2+(1-d)v$\\
$a=d$, I&$d=pu/2$&$d^2v$\\
$a=d$, II&$d=p/2+pu/2$&$2d-p+(d^2-2d+p)v$\\
$a=(t+p)/2$, I&$d=p/2+qu$&$(2d-p)v$\\
$a=(t+p)/2$, II&$d=h+pu/2$&$2d+p-2+2qv$\\
\bottomrule
\end{tabular}
\end{center}
These cubes cover every feasible boundary point.  Some contain extra
infeasible points; positivity there only strengthens the certificate.  In
region $C$, the last coordinate on the two $a=(t+p)/2$ faces is split at
$v=1/3$ to isolate the $p=2/3$ equality.

The high-density region uses coordinates adapted to its curved feasible
edges.  On $a=d$, after the first cube above, write
\begin{equation}
 s=1-d,\qquad q=s^2u,\qquad
 t=2d-p+(s^2-q)v.
 \label{eq:high-upper-d}
\end{equation}
The additional constraints $q\le1/3$ and $2d-p\ge0$ give an exact cover.
On $a=(t+p)/2$, write
\begin{equation}
 q=\frac X3,\quad p=1-q,\quad s=1-d,\quad
 t=p-2s+2qv.
 \label{eq:high-upper-t}
\end{equation}
The remaining feasibility conditions are exactly
\begin{equation}
 s-qv\ge0,\qquad s^2+q-2qv\ge0.
 \label{eq:high-upper-t-constraints}
\end{equation}
The rectangle is split at $s=q$ and $v=q$, producing $ll,lh,rl,rh$.
The lower face $a=d-q=p-s$ is similarly split at $s=q$ and uses
$t=p-2s+sv$.  These substitutions remove the square-root boundary that
would otherwise occur on $a=d$ and expose every endpoint zero in rational
coordinates.

\section{Exact Bernstein certificate}

If a polynomial $P$ on a box is written in the tensor Bernstein basis, then
\[
 P=\sum_{\boldsymbol i}b_{\boldsymbol i}
 \prod_j B_{i_j,n_j}(u_j).
\]
The basis functions are nonnegative and sum to one.  Thus nonnegative
$b_{\boldsymbol i}$ prove $P\ge0$ on the box.  Exact de Casteljau bisection
gives the coefficients on each dyadic child.  Conversely, if every
coefficient of a feasibility constraint is strictly negative, that box is
infeasible.  The verifier clears a positive common denominator and performs
de Casteljau subdivision over the integers.  A separate exact check proves
equivalence to rational subdivision up to a positive scale.

After each face substitution, the residual is multiplied by $d(1-d)$; in
the interior it is multiplied by $4dt$.  All other denominators depend only
on the region parameter, and their strict signs are proved by univariate
Bernstein coefficients.  The certificate tree is:
\begin{center}
\begin{tabular}{llrrr}
\toprule
region&piece&visited&accepted&infeasible\\
\midrule
$A$&interior&1097&261&288\\
&lower I, II&34&18&0\\
&upper-$d$ I, II&6&4&0\\
&upper-$t$ I, II&538&270&0\\
\midrule
$C$&interior&479&106&134\\
&lower I, II&36&19&0\\
&upper-$d$ I, II&6&4&0\\
&upper-$t$ I low/high&7704&3853&0\\
&upper-$t$ II low/high&7210&3606&0\\
\midrule
$H$&interior&221&42&69\\
&lower faces&37&20&0\\
&upper-$d$ faces&2&2&0\\
&upper-$t$ first range&44&23&0\\
&upper-$t$ split $ll,lh,rl,rh$&11670&5836&1\\
\midrule
&total&29084&14064&492\\
\bottomrule
\end{tabular}
\end{center}
Accepted counts leaf boxes proved directly or by an exact local lemma;
internal nodes are included in visited.

There are genuine equality boxes, so subdivision alone is not used as a
substitute for endpoint proofs.  The verifier checks:
\begin{itemize}
\item At the region-$A$ interior equality $x=0$, after
  $d=1/2+y$ and $v=1/4+z$, the quadratic part is
  \[
  5y^2-8yz+8z^2=y^2+4(y-z)^2+4z^2.
  \]
  The higher remainder is dominated on the box of radius $1/64$, while the
  quotient by $x$ has 380 nonnegative tensor Bernstein coefficients there.
\item The isolated region-$A$ upper-face corner has leading terms, after a
  positive common scale,
  \[
  112500z+86021x^2-195000xy+112500y^2.
  \]
  The quadratic matrix remains positive definite after subtracting
  $500(x^2+y^2)$; the exact remainder radius is $1/2048$.
\item Each of the four region-$C$ corners at $p=2/3$ has a positive-definite
  leading form in $((1-w)^2,u,z)$.  Sylvester's criterion is checked over
  the rationals after subtracting the margins
  $19/100,19/100,3/20,9/50$; the exact radii are
  $1/128,1/256,1/512,1/1024$.
\item On the high $ll$ piece a weighted quadratic form controls the endpoint
  uniformly in $p$.  On the high $rh$ piece there are separate exact
  estimates at $p=1$ and $p=2/3$, together with 13 rational dyadic parameter
  bands.  The remaining high upper-face corner has the positive-cone form
  \[
  \frac{119r^2+100ru-88rz+28u^2-16uz+32z^2}{1458}.
  \]
  Dropping the nonnegative $100ru$ term leaves a matrix positive definite
  after subtracting $I/200$, with exact radius $1/128$.
\end{itemize}
All polynomial construction, coefficient extraction, feasibility pruning,
local identities, margins, remainders, tree sizes, graph identification, and
target algebra are checked by
\begin{center}
\url{codes/verify_atlas126_triangle_c4_vertex_supporting_plane.py}.
\end{center}
There is no floating-point sign decision.  This proves
Lemma~\ref{lem:scalar}.

\section{Integration and endpoints}
\label{sec:endpoints}

Integrating \eqref{eq:supporting-plane} and using \eqref{eq:moments} gives
\begin{align*}
 \int F_p(d,a,t)
 &\ge C(p)
 +\beta\left(\int d-p\right)
 +\gamma\left(\int a-\int d^2\right)
 +\lambda\left(\int t-g\right)\\
 &\ge C(p).
\end{align*}
The last inequality is \eqref{eq:reiher} and $\lambda\ge0$ on the low
interval; on the high interval the last term is absent.  Combining this with
\eqref{eq:rooted-reduction} proves the desired graphon inequality once the
rooted endpoints are interpreted as follows.

If $d(x)=0$, then $W(x,\cdot)=0$ almost everywhere, so $a(x)=t(x)=0$ and
the projected contribution is $0$.  Its scalar gap is
$G_0=\beta p+\lambda g-C\ge0$.  If $d(x)=1$, then
$W(x,\cdot)=1$ almost everywhere, so $a(x)=t(x)=p$, the rooted $C_4$
projection is $p^2$, and $F=p^3$.  Its scalar gap is $G_1\ge0$.
When $t=0$ and $0<d<1$, define $F=0$; the relevant minima occur on the
certified boundary faces.  This removes every zero denominator without
appealing to continuity outside the actual feasible set.

At $p=1/2$, $C(p)=0$; region $A$ includes this endpoint and its exact local
certificate, with $g=0$ and $\lambda=1/16$.  At $p=2/3$, the region-$C$ and
region-$H$ coefficients agree as in \eqref{eq:join-coefficients}; all four
transition corners are audited.  If $p=1$, then $W=1$ almost everywhere and
both sides are $1$.

\begin{theorem}
For the triangle--$C_4$ one-vertex amalgam $H$ (Graph Atlas 126) and every
symmetric measurable graphon $W$ of edge density $p\in[1/2,1]$,
\[
 t(H,W)\ge p^2(2p-1)\bigl(p^3+(1-p)^3\bigr).
\]
\end{theorem}

\section{Audit boundary and formalization status}

The proof is graph-specific.  It uses the factorization at the common vertex
and a supporting plane tailored to three scalar rooted statistics.  It does
not assert a general vertex-amalgamation, cone, or join closure theorem.  On
the low interval, $\lambda$ is calibrated sharply to the triangle-density
profile; a coarser triangle bound does not prove the target.

The companion script also performs 46 exact two-block graphon checks.  Each
independently enumerates the homomorphism density, verifies the rooted
factorization and scalar constraints, and checks the projection and target
inequalities.  They are sanity checks, not a replacement for the
arbitrary-graphon argument.

No Lean command was run and no claim of Lean verification is made.  A future
formalization should define the graph in
\texttt{lean/Taeyoung/Examples/Graph126.lean}, import or prove the sharp
triangle profile, formalize the weighted Cauchy projection, and reflect the
exact rational scalar certificates.

\begin{thebibliography}{9}
\bibitem{ReiherCliqueDensity}
Christian Reiher,
\emph{The clique density theorem},
Annals of Mathematics \textbf{184} (2016), 683--707.
\end{thebibliography}
\end{document}
